\chapter{Experiment Results}
\label{chap:experiment-results}

\section{Main Comparison Results}
This section compares the proposed method with the baselines on TextVQA, DocVQA and ChartQA.

\subsection{Baseline Results}
Before applying any adaptation or distillation, we report the original performance of the two base student models, \textbf{SmolVLM-500M} and \textbf{SmolVLM-256M}. These baseline results establish the starting point for later comparison. All results are obtained using \textbf{VLM-Eval-Kit} under the same evaluation pipeline described in Chapter~\ref{chap:evaluation-methodology}.

\begin{table}[htbp]
\centering
\begin{tabular}{lccc}
\toprule
\textbf{Model} & \textbf{TextVQA} & \textbf{DocVQA} & \textbf{ChartQA} \\
\midrule
SmolVLM-500M & 56.75 & 59.5247 & 62.8 \\
SmolVLM-256M & 47.17 & 48.6411 & 45.04 \\
\bottomrule
\end{tabular}
\caption{Original benchmark results of the two SmolVLM models evaluated with VLM-Eval-Kit.}
\label{tab:smolvlm-original-results}
\end{table}

\section{Ablation Study}
This section analyzes the effect of teacher-set selection, weighting strategy and student model size on final performance.

\section{Qualitative Analysis}
This section presents representative examples to inspect how distillation affects text reading and basic reasoning behavior.

\section{Discussion}
This section discusses the trade-off between model size, teacher supervision and transfer performance across the evaluated benchmarks.
